\documentclass[11pt,a4paper]{article}
\usepackage[margin=2.5cm]{geometry}

\usepackage[T1]{fontenc}   
\usepackage[utf8]{inputenc} 
\usepackage[english]{babel}

\usepackage{newtxtext,newtxmath}

\usepackage{microtype}

\usepackage{xcolor}
\definecolor{citationred}{RGB}{150,0,0}

\usepackage[
  colorlinks=true,
  citecolor=citationred,
  linkcolor=black,
  urlcolor=black
]{hyperref}

\usepackage[authoryear,round]{natbib}



\usepackage[most]{tcolorbox}  
\tcbuselibrary{skins}         


\newtcolorbox{commentbox}{
  colframe=green,        
  boxrule=0.4pt,         
  sharp corners,         
  enhanced,              
  left=2mm,right=2mm,    
  bottom=2mm,
  top=1.9em,             
  overlay={
    \node[anchor=north west,fill=citationred,text=white,
          font=\bfseries,inner ysep=2pt,inner xsep=4pt]
         at (frame.north west) {Commentaires};
  },
}

\renewcommand{\thefootnote}{\fnsymbol{footnote}}

\begin{document}


\section*{Deep values: most people share some degree of universalism}



% 2. Deep values: most people share some degree of universalism

% =======================================
% - general questions (ISSP, UNDP, cf. A.1.1 de mon papier NHB)
% =======================================

ISSP surveys indicate that concern for global problems is displayed by a majority. Across 36 countries, 82\% of respondents agree that environmental problems require international agreements that countries must follow, while only 4\% disagree. In a more recent survey covering 29 countries, 78\% respondents consider existing economic differences between rich and poor countries too large, while 5\% who disagree \citep{issp_environment_2010,issp_inequality_2019}. 

% =======================================
% - wrong perceptions (Fehr et al 2022)
% =======================================
Despite such concerns, there is a gap between attitudes toward global inequalities and actual policies. A first attempt to explain it comes from misperception. Indeed, Germans underestimate their global income rank by an average of 15 percentage points. Correcting perceived national rank changes, both national and global, redistributive preferences among left-leaning respondents, whereas correcting perceived global rank has no detectable effect on support for global redistribution. In addition, 90\% of respondents expressed a preference for some degree of global redistribution, although intensity vary across individuals.
By contrast, aid support in US is driven by information on global distribution, because people underestimate their rank by 27 percentage points in their global income rank and overestimate global median income by a factor 10, they tend to increase their support to foreign aid by almost 10 percentage points and donation to international charities by 55\% \citep{nair_misperceptions_2018}.

% =======================================
% - universalism (Enke, etc.), cf. A.1.5 
% =======================================

Literature has long studied the concept of global identity (\citet{rey_2018} for a review). Universalism has several definitions in the literature, although some are extremely close. \citet{enke_moral_2023}, \citet{enke_universalism_2024}, and \citet{cappelen_universalism_2025} define it as the extent to which altruism remains invariant to recipients' identity, group membership, or social and geographic distance. \citet{boyer-kassem_how_2026} retain this conceptual definition but argue that its measurement requires recipients to be economically comparable and to belong to clearly identified groups. A related concept is the moral circle \citep{waytz_ideological_2019,jaeger_toward_2026}.

\citet{bayram_2015} reports that ``78\% of the participants in 57 countries see themselves as citizens of the world'', referring to the 2005--2008 wave of the World Values Survey. In the 2017--2022 wave, however, 84\% of respondents report feeling close or very close to their country, compared with 47\% feeling close to the world\footnote{These results were calculated using the World Values Survey's online analysis tool, available at \url{https://www.worldvaluessurvey.org/WVSOnline.jsp}.}. 

Universalism is reflected not only in stated preferences but also in observed patterns of giving. Using approximately four million donations made through the DonorsChoose platform, worth about \$305 million, \citet{enke_universalism_2024} measure universalism at the U.S. congressional-district level as the extent to which charitable giving remains invariant to geographic and social distance. This district-level measure predicts Democratic vote shares, legislators' roll-call voting, and the moral content of congressional speeches more strongly than district income or education. The individual-level allocation tasks come from a separate paper. \citet{enke_moral_2023} ask 11,000 respondents in Australia, France, Germany, Sweden, and the United States to allocate hypothetical money and trust points between equally rich recipients at different social distances. Their composite measure ranges from zero, when all money and trust are allocated to in-group members, to one, with 0.5 representing equal treatment. Average scores are below this equal-treatment benchmark in every country, ranging from 0.38 in Australia and Sweden to 0.41 in France, with intermediate values of 0.39 in the United States and 0.40 in Germany.\footnote{\citet{boyer-kassem_how_2026} argue that this design does not fully isolate social distance. Equal nominal incomes do not necessarily imply comparable purchasing power, wealth, or relative income positions across countries. Moreover, a ``random stranger'' may share one or more relevant group memberships with the respondent. Differences in allocations may therefore reflect economic disparities or ambiguously defined group boundaries in addition to universalism.}

\citet{prather_2013} controls for inequality and still finds a nationalist bias in the United States, where 84\% of Americans agree that ``taking care of problems at home is more important than giving aid to foreign countries'' \citep{PIPA2001}. Support for a U.S. program to fight hunger is 10 percentage points lower when the hypothetical program operates abroad rather than domestically, and the gap reaches 20 percentage points when the choice concerns cutting an existing program.

Using a similar approach, \citet{enke_moral_2023} show that moral universalism as the extent to which individuals exhibit the same level of altruism and trust toward strangers as toward members of their in-groups, and explains political views better than standard political economy variables such as income, wealth or perception about government efficiency across Australia, Germany, France, the US, and Sweden. For instance, if universalism levels were the same for all, individuals' policy position would not be as much polarized as they are. 
Similarly, using surveys covering 60 countries and 64,000 respondents, \citet{cappelen_universalism_2025} ask respondents to distribute the local-currency equivalent of a hypothetical \$1,000 between two recipients with the same standard of living. Each respondent completes two randomly selected domestic tasks comparing a stranger from their country with either a family member, friend, neighbour, co-religionist, or co-ethnic, followed by a foreign task comparing a domestic stranger with a stranger living elsewhere in the world. The authors define fully universalist respondents as those who divide the money equally in every task: approximately 26\% satisfy this criterion, whereas 17\% allocate, on average across the tasks, no more than 20\% to the stranger. Country averages also vary substantially: the share allocated to the more socially distant recipient is approximately 26\% in China, India, and Israel, compared with 46\% in Ethiopia \citep[Figures 3--5]{cappelen_universalism_2025}.

%LB: j'avais parlé de climat puisque c'est la politique climatique qui dans l'idée s'approche le plus de ce à quoi on s'intéresse ici. 
Universalism can be surprising because it may lead people to sacrifice their own well-being for the sake of distant groups of people. In a lab experiment with U.S. students, \citet{cherry_2017} find that 61.1\% of participants support egalitarian policies even when these reduce their own material payoff. Complementary international experimental evidence from a lab experiment matching participants from Norway and Germany with participants from Uganda and Tanzania in a dictator game shows that participants acted as moral cosmopolitans and did not assign importance to nationality in their distributive choices \citep{cappelen_needs_2013}. Likewise, \citet{leiserowitz_2006} reports that a majority of Americans express greater concern about the consequences of climate change for people worldwide (50\%) or for non-human nature (18\%) than for themselves and their families (12\%) or for the United States (9\%). Consistent with these findings, a 2017 Focus 2030 survey shows that 40\% of French respondents agree that reducing poverty in developing countries should be a priority for the European Union, compared with only 19\% who disagree.

% =======================================
%- moral circle, group defended: Waytz et al 19, Jaeger & Wilks 23 (give descriptive stats and not only correlation if possible), Fabre et al 25, Fabre 26
% =======================================

A complementary way to describe this heterogeneity is through the concept of the ``moral circle'' defined as the set of people or other entities considered worthy of moral consideration.
Across seven study samples with U.S. respondents, \citet{waytz_ideological_2019} show that left-leaning individuals consistently report broader moral circles than conservative respondents. They ask their respondents to rate, on a scale from one to five, how strongly they identify with their community, their country, and humanity as a whole. In the full sample, average identification is slightly higher with humanity (3.25) than with the country (3.09) or community (3.06). Among liberal respondents, defined as positions one to three on a seven-point ideology scale, average identification with humanity is 3.39, compared with 2.99 for the country. Among conservatives, defined as positions five to seven, the pattern reverses, with averages of 2.64 for humanity and 3.54 for the country \citep[Figure 3]{waytz_ideological_2019}. 

In two studies in the Netherlands, US, UK and Australia, \citet{jaeger_relative_2023} decompose variation in moral concern into three components. Dutch students evaluate the moral standing of 30 animal targets. Systematic differences between targets explain 40\% of the variation: some animals are consistently considered more worthy of moral concern than others. Differences between judges explain 29\%: some respondents generally assign more moral concern than others. Finally, judge-by-target interactions explain 27\%, meaning that particular respondents evaluate particular targets differently from what their general tendencies and the targets' average ratings would predict \citep[Figures 1]{jaeger_relative_2023}.

We also find this structure when respondents are asked which group they advocate for when they vote. Among 8,000 respondents in five Western countries, \citet{fabre_majority_2025} argue that ``When we ask respondents which group they defend when they vote, 20\% choose ‘sentient beings (humans and animals)’, 22\% choose ‘humans’, 33\% select their ‘fellow citizens’ (or ‘Europeans’), 15\% choose ‘My family and myself’ and the remaining 9\% choose another group (mainly ‘My State or region’ or ‘People sharing my culture or religion’). Notably, a majority of left-wing voters choose humans or sentient beings.''. In a broader survey of 12,000 respondents across eleven 
high-income 
countries, \citet{fabre_public_2026} uses a new question to measure universalism asking respondents the group they advocate when voting and ``45\% choose a universalist  response (either “Humans” or “Sentient beings (humans and animals)”), while 32\% opt for their fellow citizens''. The universalist share reaches 50\% in Europe, 57\% in Saudi Arabia, and 59\% among left-wing respondents \citep{fabre_public_2026}. Universalistic concern is widespread, despite median voters favoring their fellow citizens in the U.S. and Japan.

% =======================================
%- why HICs should help LICs: duty (Fabre 26)
% =======================================


When respondents are asked why ``high-income countries should support low-income countries'', they first invoke a moral duty, on average at 54\% at the exception of France (43\%) and Japan (45\%), while, on average, 38\% of respondents suggest the ``Long-term interest of HICs to help LICs''. Finally, only 16\% of respondents advocate non of these arguments, which puts an end to self-interest or guilt bias \citep{fabre_public_2026}.


\newpage
\bibliographystyle{plainnat}
\bibliography{chapter/references}

\end{document}
